\section{$q$-Binomial formulas}

\subsection{Combinatorial definition of $q$-binomial coefficients}

This section formalizes Definition~\ref{def.pars.qbinom.qbinom} from the combinatorial
side: the $q$-binomial as a sum over partitions, and proves that this agrees with
the recurrence-based definition from the previous chapter.

\begin{definition}
\label{def.qbinomial-formulas.qbinom-poly}
\lean{AlgebraicCombinatorics.QBinomialRec.qBinomialPolyDef, AlgebraicCombinatorics.QBinomialRec.qBinomialEval}
\leanhelper
\leanok
The $q$-binomial coefficient as a polynomial in $\mathbb{Z}[q]$ via the partition sum:
\[
\binom{n}{k}_q = \sum_{s=0}^{k(n-k)} c(s, k, n-k) \cdot q^s
\]
where $c(s,k,m)$ is the number of partitions of $s$ fitting in a $k \times m$ box.
The evaluation $\binom{n}{k}_a$ is obtained by substituting $a$ for $q$.
\end{definition}

\begin{lemma}
\label{lem.qbinomial-formulas.pascal-poly}
\lean{AlgebraicCombinatorics.QBinomialRec.qBinomialPolyDef_pascal}
\leanhelper
\leanok
The combinatorial $q$-binomial satisfies the $q$-Pascal identity as polynomials:
\[
\binom{n+1}{k+1}_q = \binom{n}{k}_q + q^{k+1} \cdot \binom{n}{k+1}_q.
\]
This is the polynomial version; the evaluated version follows by evaluation.
\end{lemma}

\begin{proof}
\leanok
Partitions in a $(k{+}1) \times (n{-}k)$ box split into those with $\le k$ parts
and those with exactly $k{+}1$ parts. The latter biject with partitions in a
$(k{+}1) \times (n{-}k{-}1)$ box via ``removing a column'' (subtracting 1 from
each of the $k{+}1$ parts), with size decreased by $k{+}1$ (giving the factor $q^{k+1}$).
\end{proof}

\begin{lemma}
\label{lem.qbinomial-formulas.exact-parts-bij}
\lean{AlgebraicCombinatorics.QBinomialRec.partitionsInBoxExact_card_eq}
\leanhelper
\leanok
For $m \ge 1$ and $s \ge k+1$, the number of partitions of $s$ with exactly $k+1$ parts,
each $\le m$, equals the number of partitions of $s - (k+1)$ fitting in a
$(k+1) \times (m-1)$ box:
\[
|\mathrm{partitionsInBoxExact}(s, k{+}1, m)| = c(s{-}(k{+}1), k{+}1, m{-}1).
\]
\end{lemma}

\begin{proof}
\leanok
The forward bijection subtracts 1 from each of the $k{+}1$ parts (and filters zeros);
the inverse adds 1 to each part and pads with 1s to reach exactly $k{+}1$ parts.
Injectivity follows from the recovery lemma for filtered multiset subtraction;
surjectivity is verified directly.
\end{proof}

\begin{theorem}
\label{thm.qbinomial-formulas.eval-eq-qbinomial}
\lean{AlgebraicCombinatorics.QBinomialRec.qBinomialEval_eq_qBinomial}
\leanhelper
\leanok
The combinatorial definition (sum over partitions) agrees with the recurrence-based
definition from Definition~\ref{def.pars.qbinom.qbinom}:
for all $n, k \in \mathbb{N}$ and ring element $q$,
\[
\mathrm{qBinomialEval}(n, k, q) = \binom{n}{k}_q.
\]
\end{theorem}

\begin{proof}
\leanok
Both definitions satisfy the same boundary conditions ($\binom{n}{0}_q = 1$,
$\binom{n}{k}_q = 0$ for $k > n$) and the same $q$-Pascal recurrence.
By strong induction on $n$ with case analysis on $k$, the two definitions agree.
\end{proof}

\subsection{Alternative characterizations from the recurrence}

\begin{lemma}
\label{lem.qbinomial-formulas.sum-monotone}
\lean{AlgebraicCombinatorics.QBinomialRec.qBinomial_eq_sum_monotone}
\leanhelper
\leanok
For $k \le n$, the recurrence-based $q$-binomial equals the sum over monotone functions:
\[
\binom{n}{k}_q = \sum_{f \in \mathrm{monotoneFunctions}(k, n{-}k)} q^{\mathrm{sumValues}(f)}.
\]
\end{lemma}

\begin{proof}
\leanok
By induction on $n$, splitting the monotone functions into those starting at 0
and those starting at a positive value, using the bijections
$\mathrm{dropFirst}$/$\mathrm{prependZero}$ and $\mathrm{shiftDown}$/$\mathrm{shiftUp}$.
\end{proof}

\begin{lemma}
\label{lem.qbinomial-formulas.sum-subsets}
\lean{AlgebraicCombinatorics.QBinomialRec.qBinomial_eq_sum_subsets}
\leanhelper
\leanok
For all $n, k$:
\[
\binom{n}{k}_q = \sum_{\substack{S \subseteq \{1,\ldots,n\} \\ |S|=k}}
q^{\mathrm{sum}(S) - (1+2+\cdots+k)}.
\]
\end{lemma}

\begin{proof}
\leanok
Via the bijection from monotone $k$-tuples in $\{0,\ldots,n{-}k\}$ to $k$-element
subsets of $\{1,\ldots,n\}$: $(i_1,\ldots,i_k) \mapsto \{i_1{+}1, i_2{+}2, \ldots, i_k{+}k\}$.
For $k > n$, both sides are 0.
\end{proof}

\begin{lemma}
\label{lem.qbinomial-formulas.prod-div}
\lean{AlgebraicCombinatorics.QBinomialRec.qBinomial_eq_prod_div}
\leanhelper
\leanok
Over a field, when $[i{+}1]_q \ne 0$ for $i = 0, \ldots, k{-}1$:
\[
\binom{n}{k}_q = \prod_{i=0}^{k-1} \frac{[n-i]_q}{[i+1]_q}.
\]
\end{lemma}

\begin{proof}
\leanok
By strong induction on $n$, using the $q$-Pascal recurrence and the splitting identity
$[n{+}1]_q = [k{+}1]_q + q^{k+1} [n{-}k]_q$.
\end{proof}

\subsection{Product expansion rule}

\begin{lemma}
\label{lem.prodrule.sum-ai-plus-bi}
\lean{AlgebraicCombinatorics.prod_add_eq_sum_over_subsets}
\leantarget
\leanok
Let $L$ be a commutative ring. Let
$n\in\mathbb{N}$. Let $\left[  n\right]  $ denote the set $\left\{
1,2,\ldots,n\right\}  $. Let $a_{1},a_{2},\ldots,a_{n}$ be $n$ elements of
$L$. Let $b_{1},b_{2},\ldots,b_{n}$ be $n$ further elements of $L$. Then,
\begin{equation}
\prod_{i=1}^{n}\left(  a_{i}+b_{i}\right)  =\sum_{S\subseteq\left[  n\right]
}\left(  \prod_{i\in S}a_{i}\right)  \left(  \prod_{i\in\left[  n\right]
\setminus S}b_{i}\right)  . \label{eq.lem.prodrule.sum-ai-plus-bi.eq}
\end{equation}
\end{lemma}

\begin{proof}
\leanok
When expanding the product $\prod_{i=1}^{n}\left(  a_{i}+b_{i}\right)
=\left(  a_{1}+b_{1}\right)  \left(  a_{2}+b_{2}\right)  \cdots\left(
a_{n}+b_{n}\right)  $, you obtain a sum of $2^{n}$ terms, each of which is a
product of one addend chosen from each of the $n$ sums $a_{1}+b_{1}%
,a_{2}+b_{2},\ldots,a_{n}+b_{n}$. This is precisely what the right hand side
of \eqref{eq.lem.prodrule.sum-ai-plus-bi.eq} is.
A rigorous proof can be found in \cite[Exercise 6.1 \textbf{(a)}]{detnotes}.
The formal proof follows by induction on $n$, using the distributive law at each step.
\end{proof}

\subsection{First $q$-binomial theorem}

\subsubsection{Helpers for Theorem~\ref{thm.pars.qbinom.binom1}}

\begin{lemma}
\label{lem.qbinomial-formulas.subset-sum-formula}
\lean{AlgebraicCombinatorics.QBinomialRec.qBinomial_sum_subsets}
\leanhelper
\leanok
For all $n, k \in \mathbb{N}$:
\[
\sum_{\substack{S \subseteq \{0,1,\ldots,n-1\} \\ |S|=k}} q^{\sum_{i \in S} i}
= q^{k(k-1)/2} \cdot \binom{n}{k}_q.
\]
This relates the subset sum over $\{0,1,\ldots,n-1\}$ to the $q$-binomial coefficient.
\end{lemma}

\begin{proof}
\leanok
By induction on $n$, using the $q$-Pascal recurrence. The subsets of $\{0,\ldots,n\}$
of size $k{+}1$ split into those containing $n$ and those not containing $n$.
The key algebraic identity is
$q^k \binom{n}{k+1}_q + q^n \binom{n}{k}_q = q^k \binom{n}{k}_q + q^{k+1} \cdot q^k \binom{n}{k+1}_q$,
which is proved by induction using algebraic manipulation.
\end{proof}

\begin{lemma}
\label{lem.qbinomial-formulas.qbinom-induction-identity}
\lean{AlgebraicCombinatorics.QBinomialRec.qBinomial_induction_identity}
\leanhelper
\leanok
For all $n, k \in \mathbb{N}$:
\[
q^k \binom{n}{k+1}_q + q^n \binom{n}{k}_q
= q^k \binom{n}{k}_q + q^{k+1} \cdot q^k \binom{n}{k+1}_q.
\]
\end{lemma}

\begin{proof}
\leanok
By induction on $n$, using the $q$-Pascal recurrence and algebraic manipulation.
\end{proof}

\begin{theorem}
[1st $q$-binomial theorem]
\label{thm.pars.qbinom.binom1}
\lean{AlgebraicCombinatorics.qBinomial_first_theorem}
\leantarget
\leanok
Let $K$ be a
commutative ring. Let $a,b\in K$ and $n\in\mathbb{N}$. In the polynomial ring
$K\left[  q\right]  $, we have
\[
\left(  aq^{0}+b\right)  \left(  aq^{1}+b\right)  \cdots\left(  aq^{n-1}%
+b\right)  =\sum_{k=0}^{n}q^{k\left(  k-1\right)  /2}\binom{n}{k}_{q}%
a^{k}b^{n-k}.
\]
\end{theorem}

\begin{proof}
\leanok
Let $\left[  n\right]  $ denote the set $\left\{  1,2,\ldots,n\right\}  $. We have
\begin{align*}
&  \left(  aq^{0}+b\right)  \left(  aq^{1}+b\right)  \cdots\left(
aq^{n-1}+b\right)  =\prod_{i=1}^{n}\left(  aq^{i-1}+b\right) \\
&  =\sum_{S\subseteq\left[  n\right]  }\left(  \prod_{i\in
S}\left(  aq^{i-1}\right)  \right) \left(  \prod_{i\in\left[  n\right]
\setminus S}b\right)
\end{align*}
by Lemma~\ref{lem.prodrule.sum-ai-plus-bi}, applied to $L=K\left[  q\right]$,
$a_{i}=aq^{i-1}$ and $b_{i}=b$.
Now,
\begin{align*}
&  =\sum_{S\subseteq\left[  n\right]  }
a^{\left\vert S\right\vert }
\left(  \prod_{i\in S}q^{i-1}\right)
b^{\left\vert \left[  n\right]  \setminus S\right\vert }
=\sum_{k=0}^{n}\ \ \sum_{\substack{S\subseteq\left[  n\right]
;\\\left\vert S\right\vert =k}}
a^{k}
q^{\operatorname{sum}S-k}
b^{n-k}
\end{align*}
(since $\prod_{i\in S}q^{i-1} = q^{\operatorname{sum}S - |S|}$
and $|[n]\setminus S| = n - k$).
Using
$q^{\operatorname{sum}S-k}
= q^{\operatorname{sum}S-(1+2+\cdots+k)} \cdot q^{k(k-1)/2}$
(since $1+2+\cdots+(k-1) = k(k-1)/2$),
we obtain
\begin{align*}
&  =\sum_{k=0}^{n}q^{k\left(  k-1\right)  /2}\underbrace{\left(
\sum\limits_{\substack{S\subseteq\left[  n\right]  ;\\\left\vert S\right\vert
=k}}q^{\operatorname{sum}S-\left(  1+2+\cdots+k\right)  }\right)
}_{=\binom{n}{k}_{q}}a^{k}b^{n-k}
=\sum_{k=0}^{n}q^{k\left(  k-1\right)  /2}\binom{n}{k}_{q}a^{k}b^{n-k}.
\end{align*}
(Here, the underbrace equality follows from Proposition
\ref{prop.pars.qbinom.alt-defs} \textbf{(b)}.)
This proves Theorem~\ref{thm.pars.qbinom.binom1}.
\end{proof}

\begin{lemma}
\label{lem.qbinomial-formulas.binom1-at-one}
\lean{AlgebraicCombinatorics.QBinomialRec.qBinomial_first_theorem_at_one}
\leanhelper
\leanok
The first $q$-binomial theorem at $q = 1$ recovers the ordinary binomial formula:
\[
(a+b)^n = \sum_{k=0}^{n} \binom{n}{k} a^k b^{n-k}.
\]
\end{lemma}

\begin{proof}
\leanok
Specializing Theorem~\ref{thm.pars.qbinom.binom1} at $q = 1$: the LHS becomes
$(a+b)^n$ and each $\binom{n}{k}_1 = \binom{n}{k}$.
\end{proof}

\subsection{Second $q$-binomial theorem}

\subsubsection{Helpers for Theorem~\ref{thm.pars.qbinom.binom2}}

\begin{lemma}
\label{lem.qbinomial-formulas.b-mul-pow-a}
\lean{AlgebraicCombinatorics.QBinomialRec.b_mul_pow_a}
\leanhelper
\leanok
Let $A$ be an $L$-algebra, let $\omega \in L$, and let $a, b \in A$ satisfy $ba = \omega \cdot ab$.
Then for all $k \in \mathbb{N}$:
\[
b \cdot a^k = \omega^k \cdot (a^k \cdot b).
\]
\end{lemma}

\begin{proof}
\leanok
By induction on $k$. The base case is trivial. For the inductive step:
$b \cdot a^{k+1} = b \cdot a^k \cdot a = \omega^k (a^k b) \cdot a
= \omega^k \cdot a^k \cdot (ba) = \omega^k \cdot a^k \cdot \omega(ab)
= \omega^{k+1}(a^{k+1} b)$.
\end{proof}

\begin{theorem}
[2nd $q$-binomial theorem, aka Potter's binomial theorem]
\label{thm.pars.qbinom.binom2}
\lean{AlgebraicCombinatorics.qBinomial_second_theorem}
\leantarget
\leanok
Let $L$ be a commutative ring. Let $\omega\in
L$. Let $A$ be a noncommutative $L$-algebra. Let $a,b\in A$ be such that
$ba=\omega ab$. Then,
\[
\left(  a+b\right)  ^{n}=\sum_{k=0}^{n}\binom{n}{k}_{\omega}a^{k}b^{n-k}.
\]
\end{theorem}

\begin{proof}
\leanok
By induction on $n$.
The base case $n=0$ is trivial.
For the inductive step, write $(a+b)^{n+1} = (a+b) \cdot (a+b)^n$
and apply the induction hypothesis to $(a+b)^n$.
Distribute $(a+b)$ over the sum.
The key step is the relation $b \cdot a^k = \omega^k \cdot a^k \cdot b$,
which follows from $ba = \omega ab$ by induction on $k$
(proved in Lemma~\ref{lem.qbinomial-formulas.b-mul-pow-a}).
After reindexing and combining terms, the coefficients match via
the $q$-Pascal recurrence $\binom{n+1}{k+1}_\omega = \binom{n}{k}_\omega + \omega^{k+1}\binom{n}{k+1}_\omega$.
\end{proof}

\subsection{Counting subspaces of vector spaces}

\subsubsection{Helpers for Theorem~\ref{thm.pars.qbinom.subsp-count}}

\begin{definition}
\label{def.qbinomial-formulas.span-map}
\lean{AlgebraicCombinatorics.QBinomialRec.spanMap}
\leanhelper
\leanok
The \emph{span map} sends a linearly independent $k$-tuple $(v_1, \ldots, v_k)$
in $V$ to its span $\mathrm{span}(v_1, \ldots, v_k)$, viewed as a $k$-dimensional subspace.
\end{definition}

\begin{lemma}
\label{lem.qbinomial-formulas.span-fiber-card}
\lean{AlgebraicCombinatorics.QBinomialRec.spanMap_fiber_card}
\leanhelper
\leanok
Each $k$-dimensional subspace $W$ has exactly
$\prod_{i=0}^{k-1} (|F|^k - |F|^i)$ preimages under the span map.
These are precisely the ordered bases of $W$: the linearly independent
$k$-tuples in $W$ that automatically span $W$ by dimension count.
\end{lemma}

\begin{proof}
\leanok
The fiber over $W$ is in bijection with the set of linearly independent
$k$-tuples in $W$. Since $W$ is $k$-dimensional, applying
Lemma~\ref{lem.pars.qbinom.lin-ind-count} (with $W$ in place of $V$,
and $k$ in place of $n$) gives the count
$\prod_{i=0}^{k-1}(|F|^k - |F|^i)$.
\end{proof}

\begin{lemma}
\label{lem.linalg.lin-ind-via-span}
\lean{AlgebraicCombinatorics.linearIndependent_iff_not_mem_span_of_lt}
\leantarget
\leanok
Let $F$ be a field. Let $V$ be an
$F$-vector space. Let $\left(  v_{1},v_{2},\ldots,v_{k}\right)  $ be a
$k$-tuple of vectors in $V$. Then, $\left(  v_{1},v_{2},\ldots,v_{k}\right)  $
is linearly independent if and only if each $i\in\left\{  1,2,\ldots
,k\right\}  $ satisfies $v_{i}\notin\operatorname{span}\left(  v_{1}%
,v_{2},\ldots,v_{i-1}\right)  $ (where the span $\operatorname{span}\left(
{}\right)  $ of an empty list is understood to be the set $\left\{
\mathbf{0}\right\}  $ consisting only of the zero vector $\mathbf{0}$). In
other words, $\left(  v_{1},v_{2},\ldots,v_{k}\right)  $ is linearly
independent if and only if we have
\begin{align*}
v_{1}  &  \notin\underbrace{\operatorname{span}\left(  {}\right)
}_{=\left\{  \mathbf{0}\right\}  }\ \ \ \ \ \ \ \ \ \ \text{and}\\
v_{2}  &  \notin\operatorname{span}\left(  v_{1}\right)
\ \ \ \ \ \ \ \ \ \ \text{and}\\
v_{3}  &  \notin\operatorname{span}\left(  v_{1},v_{2}\right)
\ \ \ \ \ \ \ \ \ \ \text{and}\\
&  \cdots\ \ \ \ \ \ \ \ \ \ \text{and}\\
v_{k}  &  \notin\operatorname{span}\left(  v_{1},v_{2},\ldots,v_{k-1}\right)
.
\end{align*}
\end{lemma}

\begin{proof}
\leanok
We prove both directions:

$\Longrightarrow:$ Assume that the $k$-tuple $\left(  v_{1}%
,v_{2},\ldots,v_{k}\right)  $ is linearly independent. Let $i\in\left\{
1,2,\ldots,k\right\}  $. If we had $v_{i}\in\operatorname{span}\left(
v_{1},v_{2},\ldots,v_{i-1}\right)  $, then we could write $v_{i}$ in the form
$v_{i}=\alpha_{1}v_{1}+\alpha_{2}v_{2}+\cdots+\alpha_{i-1}v_{i-1}$ for some
coefficients $\alpha_{1},\alpha_{2},\ldots,\alpha_{i-1}\in F$, and therefore
$\alpha_{1}v_{1}+\cdots+\alpha_{i-1}v_{i-1}+(-1)v_{i}+0v_{i+1}+\cdots+0v_{k}=\mathbf{0}$,
which would be a nontrivial linear dependence relation
(since $v_{i}$ appears with the nonzero coefficient $-1$);
this would contradict the linear
independence of $\left(  v_{1},v_{2},\ldots,v_{k}\right)  $. Hence,
$v_{i}\notin\operatorname{span}\left(  v_{1},v_{2},\ldots,v_{i-1}\right)$.

$\Longleftarrow:$ Assume that each $i\in\left\{  1,2,\ldots,k\right\}
$ satisfies $v_{i}\notin\operatorname{span}\left(  v_{1},v_{2},\ldots
,v_{i-1}\right)  $. Assume for contradiction that the $k$-tuple $\left(  v_{1}%
,v_{2},\ldots,v_{k}\right)  $ is linearly dependent. Thus, there
exist coefficients $\beta_{1},\beta_{2},\ldots,\beta_{k}\in F$ that satisfy
$\beta_{1}v_{1}+\beta_{2}v_{2}+\cdots+\beta_{k}v_{k}=\mathbf{0}$ and that are
not all zero.
Pick the \textbf{largest} $i$ with $\beta_{i}\neq0$. Then $\beta_{i+1}=\beta_{i+2}=\cdots=\beta_{k}=0$, so
$\beta_{i}v_{i} = -\left(  \beta_{1}v_{1}+\beta_{2}v_{2}+\cdots+\beta_{i-1}v_{i-1}\right)
\in\operatorname{span}\left(  v_{1},v_{2},\ldots,v_{i-1}\right)$.
Since $\beta_{i}\neq0$, we obtain $v_{i}\in\operatorname{span}\left(
v_{1},v_{2},\ldots,v_{i-1}\right)  $, contradicting our assumption.
\end{proof}

\begin{lemma}
\label{lem.pars.qbinom.lin-ind-count}
\lean{AlgebraicCombinatorics.card_linearIndependent_tuples}
\leantarget
\leanok
Let $F$ be a finite field. Let
$n,k\in\mathbb{N}$. Let $V$ be an $n$-dimensional $F$-vector space. Then,
\begin{align*}
&  \left(  \text{\# of linearly independent }k\text{-tuples of vectors in
}V\right) \\
&  =\left(  \left\vert F\right\vert ^{n}-\left\vert F\right\vert ^{0}\right)
\left(  \left\vert F\right\vert ^{n}-\left\vert F\right\vert ^{1}\right)
\cdots\left(  \left\vert F\right\vert ^{n}-\left\vert F\right\vert
^{k-1}\right)  =\prod_{i=0}^{k-1}\left(  \left\vert F\right\vert
^{n}-\left\vert F\right\vert ^{i}\right)  .
\end{align*}
\end{lemma}

\begin{proof}
\leanok
We have $\left\vert V\right\vert =\left\vert F\right\vert ^{n}$
(since $V$ is an $n$-dimensional $F$-vector space).

By Lemma~\ref{lem.linalg.lin-ind-via-span}, a $k$-tuple $\left(
v_{1},v_{2},\ldots,v_{k}\right)  $ of vectors in $V$ is linearly independent
if and only if each $v_i \notin \operatorname{span}(v_1, \ldots, v_{i-1})$.

We construct such a tuple entry by entry:
\begin{itemize}
\item First, we choose $v_{1} \in V\setminus
\left\{  \mathbf{0}\right\}$; thus, there
are $\left\vert V\right\vert - 1$ options.
Once $v_{1}$ has been chosen, $\operatorname{span}%
\left(  v_{1}\right)  $ has dimension $1$ and thus
size $\left\vert F\right\vert ^{1}$.

\item Next, we choose $v_{2} \in V\setminus
\operatorname{span}\left(  v_{1}\right)$; thus, there are
$\left\vert V\right\vert -\left\vert F\right\vert ^{1}$ options.
Once $v_{2}$ has been chosen, $\operatorname{span}\left(
v_{1},v_{2}\right)  $ has dimension $2$ and thus size
$\left\vert F\right\vert ^{2}$.

\item And so on, until the last vector $v_{k}$ has been chosen.
\end{itemize}

The total count is
$\prod_{i=0}^{k-1}\left(  \left\vert V\right\vert -\left\vert F\right\vert
^{i}\right)  =\prod_{i=0}^{k-1}\left(  \left\vert F\right\vert ^{n}-\left\vert
F\right\vert ^{i}\right)$
(since $\left\vert V\right\vert =\left\vert F\right\vert ^{n}$).
When $k > n$, both sides are $0$ since there are no linearly independent $k$-tuples
in an $n$-dimensional space.
\end{proof}

\begin{lemma}
[Multijection principle]
\label{lem.count.multijection}
\lean{AlgebraicCombinatorics.card_eq_mul_of_fibers}
\leantarget
\leanok
Let $A$ and $B$ be two
finite sets. Let $m\in\mathbb{N}$. Let $f:A\rightarrow B$ be any map. Assume
that each $b\in B$ has exactly $m$ preimages under $f$ (that is, for each
$b\in B$, there are exactly $m$ many elements $a\in A$ such that $f\left(
a\right)  =b$). Then,
\[
\left\vert A\right\vert =m\cdot\left\vert B\right\vert .
\]
\end{lemma}

\begin{proof}
\leanok
By the fiber decomposition: $|A| = \sum_{b \in B} |f^{-1}(b)|$,
where each fiber has cardinality $m$ by assumption.
Hence $|A| = \sum_{b \in B} m = m \cdot |B|$.
\end{proof}

\begin{theorem}
\label{thm.pars.qbinom.subsp-count}
\lean{AlgebraicCombinatorics.qBinomial_subspace_count}
\leantarget
\leanok
Let $F$ be a finite field. Let
$n,k\in\mathbb{N}$. Let $V$ be an $n$-dimensional $F$-vector space. Then,
\[
\binom{n}{k}_{\left\vert F\right\vert }=\left(  \text{\# of }%
k\text{-dimensional vector subspaces of }V\right)  .
\]
\end{theorem}

\begin{proof}
\leanok
If $k>n$, then $\binom{n}{k}_{\left\vert F\right\vert }=0$ (by
Proposition~\ref{prop.pars.qbinom.0}) and $\left(  \text{\# of }%
k\text{-dimensional vector subspaces of }V\right)  =0$ (since the dimension of
a subspace of $V$ is never larger than the dimension of $V$). Hence, for the rest of
this proof, we WLOG assume that $k\leq n$.

Define a map
\begin{align*}
f:\left\{  \text{linind }k\text{-tuples of vectors in }V\right\}   &
\rightarrow\left\{  k\text{-dimensional vector subspaces of }V\right\}  ,\\
\left(  v_{1},v_{2},\ldots,v_{k}\right)   &  \mapsto\operatorname{span}%
\left(  v_{1},v_{2},\ldots,v_{k}\right)  .
\end{align*}

\textbf{Observation 1:} Each $k$-dimensional vector subspace of $V$ has
exactly
$\left(  \left\vert F\right\vert ^{k}-\left\vert F\right\vert ^{0}\right)
\left(  \left\vert F\right\vert ^{k}-\left\vert F\right\vert ^{1}\right)
\cdots\left(  \left\vert F\right\vert ^{k}-\left\vert F\right\vert
^{k-1}\right)$
preimages under $f$.

Indeed, let $W$ be a $k$-dimensional vector subspace of $V$.
A preimage of $W$ under $f$ is a linind $k$-tuple
of vectors in $W$ that spans $W$.
Since $W$ is $k$-dimensional, a $k$-tuple of vectors in $W$ spans $W$ if and only if
it is linind. Hence, the number of preimages is the number of linind $k$-tuples in $W$,
which equals
$\prod_{i=0}^{k-1}(|F|^k - |F|^i)$
by Lemma~\ref{lem.pars.qbinom.lin-ind-count} (applied to $W$ and $k$ instead
of $V$ and $n$).

By the multijection principle (Lemma~\ref{lem.count.multijection}):
\begin{align*}
&  \left(  \left\vert F\right\vert ^{k}-\left\vert F\right\vert ^{0}\right)
\left(  \left\vert F\right\vert ^{k}-\left\vert F\right\vert ^{1}\right)
\cdots\left(  \left\vert F\right\vert ^{k}-\left\vert F\right\vert
^{k-1}\right) \\
&  \qquad\cdot\left(  \text{\# of }k\text{-dimensional vector
subspaces of }V\right) \\
&  =\left(  \left\vert F\right\vert ^{n}-\left\vert F\right\vert ^{0}\right)
\left(  \left\vert F\right\vert ^{n}-\left\vert F\right\vert ^{1}\right)
\cdots\left(  \left\vert F\right\vert ^{n}-\left\vert F\right\vert
^{k-1}\right)
\end{align*}
(where the right hand side comes from Lemma~\ref{lem.pars.qbinom.lin-ind-count}).
Therefore,
\begin{align*}
&  \left(  \text{\# of }k\text{-dimensional vector subspaces of }V\right) \\
&  =\dfrac{\prod_{i=0}^{k-1}\left(  \left\vert F\right\vert ^{n}-\left\vert
F\right\vert ^{i}\right)  }{\prod_{i=0}^{k-1}\left(  \left\vert F\right\vert
^{k}-\left\vert F\right\vert ^{i}\right)  }
=\prod_{i=0}^{k-1}
\dfrac{\left\vert F\right\vert ^{n-i}-1}{\left\vert F\right\vert ^{k-i}-1}
=\dfrac{\left(  1-\left\vert F\right\vert ^{n}\right)  \left(  1-\left\vert
F\right\vert ^{n-1}\right)  \cdots\left(  1-\left\vert F\right\vert
^{n-k+1}\right)  }{\left(  1-\left\vert F\right\vert ^{k}\right)  \left(
1-\left\vert F\right\vert ^{k-1}\right)  \cdots\left(  1-\left\vert
F\right\vert ^{1}\right)  }\\
&  =\binom{n}{k}_{\left\vert F\right\vert }
\end{align*}
(by Theorem~\ref{thm.pars.qbinom.quot1} \textbf{(b)}).
\end{proof}

\subsection{Limits of $q$-binomial coefficients}

\subsubsection{Helpers for Proposition~\ref{prop.pars.qbinom.lim1}}

\begin{lemma}
\label{lem.qbinomial-formulas.coeff-stability}
\lean{AlgebraicCombinatorics.QBinomialRec.qBinomial_coeff_eq}
\leanhelper
\leanok
For all $d, k, n, m \in \mathbb{N}$ with $n \ge d + k$ and $m \ge d + k$,
the $d$-th coefficient of $\binom{n}{k}_q$ equals the $d$-th coefficient of $\binom{m}{k}_q$
(as power series).
In other words, the coefficient $[q^d]\binom{n}{k}_q$ stabilizes for $n \ge d+k$.
\end{lemma}

\begin{proof}
\leanok
By strong induction on $d + k$, using the $q$-Pascal recurrence to relate
coefficients of $\binom{n}{k}_q$ and $\binom{m}{k}_q$ through $\binom{n-1}{\cdot}_q$
and $\binom{m-1}{\cdot}_q$ at lower coefficients.
\end{proof}

\begin{lemma}
\label{lem.qbinomial-formulas.stable-eq-prod}
\lean{AlgebraicCombinatorics.QBinomialRec.qBinomial_stable_eq_prod}
\leanhelper
\leanok
The stable value of the $d$-th coefficient equals the $d$-th coefficient of
the product formula:
\[
[q^d]\binom{d+k}{k}_q = [q^d]\prod_{i=1}^{k} \frac{1}{1 - q^i}.
\]
\end{lemma}

\begin{proof}
\leanok
By induction on $d + k$. The product formula satisfies a recurrence matching
the $q$-Pascal identity: extracting the last factor $(1 - q^{k+1})^{-1}$ gives
$\prod_{i=1}^{k+1}(1-q^i)^{-1} = \prod_{i=1}^{k}(1-q^i)^{-1} + q^{k+1} \cdot \prod_{i=1}^{k+1}(1-q^i)^{-1}$.
Comparing coefficients and applying the inductive hypothesis completes the proof.
\end{proof}

\begin{lemma}
\label{lem.qbinomial-formulas.conjugation-involution}
\lean{AlgebraicCombinatorics.QBinomialRec.conjugatePartition_involution'}
\leanhelper
\leanok
Young diagram conjugation (transposition) is an involution on partitions of $n$
that interchanges ``all parts $\le k$'' with ``at most $k$ parts''.
Specifically, if $p$ is a partition of $n$ with all parts $\le k$,
then its conjugate has at most $k$ parts, and vice versa.
Moreover, conjugating twice recovers the original partition.
\end{lemma}

\begin{proof}
\leanok
The conjugate is constructed using Young diagram transposition.
The involution property follows from the fact that transposing twice recovers
the original diagram.
The interchange of ``max part'' and ``number of parts'' follows from
the relationship between row lengths and column lengths under transposition.
\end{proof}

\begin{lemma}
\label{lem.qbinomial-formulas.restricted-eq-atmost}
\lean{AlgebraicCombinatorics.QBinomialRec.restricted_card_eq_atMostKParts}
\leanhelper
\leanok
The number of partitions of $n$ with all parts $\le k$ equals the number of
partitions of $n$ with at most $k$ parts:
\[
|\{p \vdash n : \text{all parts} \le k\}| = |\{p \vdash n : \text{at most } k \text{ parts}\}|.
\]
\end{lemma}

\begin{proof}
\leanok
Young diagram conjugation provides a bijection between these two sets
(Lemma~\ref{lem.qbinomial-formulas.conjugation-involution}).
\end{proof}

\begin{theorem}
\label{thm.qbinomial-formulas.limit-eq-partition-gf}
\lean{AlgebraicCombinatorics.QBinomialRec.qBinomial_limit_eq_partition_gf}
\leanhelper
\leanok
The product formula equals the generating function for partitions with at most $k$ parts:
\[
\prod_{i=1}^{k} \frac{1}{1-q^i}
= \sum_{n=0}^{\infty} |\{p \vdash n : \text{at most } k \text{ parts}\}| \cdot q^n.
\]
\end{theorem}

\begin{proof}
\leanok
The product $\prod_{i=1}^{k}(1-q^i)^{-1}$ equals the generating function
$\sum_n |\{\text{partitions of } n \text{ with all parts } \le k\}| \cdot q^n$
by the theorem on the product form of restricted partition generating functions.
Applying Lemma~\ref{lem.qbinomial-formulas.restricted-eq-atmost} converts the restricted
partition count to the count of partitions with at most $k$ parts.
\end{proof}

\begin{proposition}
\label{prop.pars.qbinom.lim1}
\lean{AlgebraicCombinatorics.qBinomial_limit}
\leantarget
\leanok
Let $k\in\mathbb{N}$ be fixed. Then,
\[
\lim_{n\rightarrow\infty}\binom{n}{k}_{q}=\sum_{n\in\mathbb{N}}\left(
p_{0}\left(  n\right)  +p_{1}\left(  n\right)  +\cdots+p_{k}\left(  n\right)
\right)  q^{n}=\prod_{i=1}^{k}\frac{1}{1-q^{i}}.
\]
(See Definition~\ref{def.pars.pn-pkn} \textbf{(a)} for the meaning of
$p_{i}\left(  n\right)  $.)
\end{proposition}

\begin{proof}[First proof (sketched)]
\leanok
For each integer $n\geq k$, we have
\begin{align*}
\binom{n}{k}_{q}  &  =\frac{\left(  1-q^{n}\right)  \left(  1-q^{n-1}%
\right)  \cdots\left(  1-q^{n-k+1}\right)  }{\left(  1-q^{k}\right)  \left(
1-q^{k-1}\right)  \cdots\left(  1-q^{1}\right)  }
\\
&  =\frac{1}{\prod_{i=1}^{k}\left(  1-q^{i}\right)  }\cdot
\prod_{i=1}^{k}\left(  1-q^{n-k+i}\right)  .
\end{align*}
(Here, the first equality holds by Theorem~\ref{thm.pars.qbinom.quot1} \textbf{(b)}.)
Since $\lim_{n\to\infty} q^n = 0$ (coefficientwise),
for each $i \in \{1,\ldots,k\}$ we have
$\lim_{n\to\infty}(1-q^{n-k+i}) = 1$.
Hence, by Corollary~\ref{cor.fps.lim.sum-prod-k},
\[
\lim_{n\to\infty}\binom{n}{k}_q
= \frac{1}{\prod_{i=1}^k(1-q^i)} \cdot 1
= \prod_{i=1}^k \frac{1}{1-q^i}.
\]
The equality
$\sum_{n\in\mathbb{N}}(p_0(n)+p_1(n)+\cdots+p_k(n))\,q^n = \prod_{i=1}^k \frac{1}{1-q^i}$
follows from Theorem~\ref{thm.pars.main-gf-0n}
(with the letters $x$, $m$ and $k$ renamed as $q$, $k$ and $i$).

The proof establishes coefficientwise stabilization via
Lemma~\ref{lem.qbinomial-formulas.coeff-stability} (showing that for $n \ge d + k$, the $d$-th coefficient
of $\binom{n}{k}_q$ is constant) and Lemma~\ref{lem.qbinomial-formulas.stable-eq-prod}
(identifying the stable value with the product formula).
The equality with the partition generating function is proved in
Theorem~\ref{thm.qbinomial-formulas.limit-eq-partition-gf} using Young diagram conjugation
to biject partitions with parts $\le k$ and partitions with $\le k$ parts.
\end{proof}
